% Generated draft: requires mathematical and source review.
\documentclass[11pt]{article}
\usepackage[margin=25mm]{geometry}
\usepackage{fontspec}
\setmainfont{Latin Modern Roman}
\usepackage{amsmath,amssymb,mathtools,unicode-math}
\setmathfont{Latin Modern Math}
\usepackage{enumitem,needspace,xurl,hyperref}
\setlength{\parindent}{0pt}
\setlength{\parskip}{5pt}
\setlength{\emergencystretch}{3em}
\newcommand{\N}{\mathbb{N}}
\newcommand{\Z}{\mathbb{Z}}
\newcommand{\Q}{\mathbb{Q}}
\newcommand{\R}{\mathbb{R}}
\newcommand{\C}{\mathbb{C}}
\newcommand{\F}{\mathbb{F}}
\newcommand{\A}{\mathbb{A}}
\newcommand{\PP}{\mathbb{P}}
\newcommand{\E}{\mathbb{E}}
\newcommand{\eps}{\varepsilon}
\newcommand{\dd}{\,\mathrm{d}}
\begin{document}
\section*{UnsolvedMath ID: 700159}
\textbf{Source problem code:} catalogue entry
\subsection{Definitions and mathematical statement}
A variety of groups \(\mathcal V\) is the class of all groups satisfying a specified set of group identities. A group \(F\) is relatively free in \(\mathcal V\) on a set \(X\) if \(F\in\mathcal V\) has a generating map \(\iota:X\to F\) such that every map from \(X\) to a group \(G\in\mathcal V\) extends uniquely to a homomorphism \(F\to G\).

A group \(H\) is \emph{radicable} if every element has an \(n\)-th root in \(H\) for every positive integer \(n\):
\[
\forall h\in H\;\forall n\in\Z_{\geq 1}\;\exists u\in H:\quad u^n=h.
\]
In particular, roots must lie in \(H\), rather than merely in an ambient group.

For a set \(W\) of group words, the associated verbal subgroup of \(F\) is
\[
W(F)=\left\langle w(g_1,\ldots,g_{k(w)}):w\in W,\ g_1,\ldots,g_{k(w)}\in F\right\rangle,
\]
where \(k(w)\) is the number of variables of the word \(w\).

\textbf{Question (N. R. Brumberg, Kourovka Notebook Problem 3.5).} Can a subgroup of a relatively free group be radicable? In particular, can a verbal subgroup of a relatively free group be radicable? With unrestricted existential quantifiers, these ask whether there exist a variety \(\mathcal V\), a set \(X\), and a relatively free group \(F\) in \(\mathcal V\) on \(X\) satisfying, respectively,
\[
\exists H\leq F\;\forall h\in H\;\forall n\in\Z_{\geq1}\;\exists u\in H:\quad u^n=h,
\]
or
\[
\exists W\;\forall h\in W(F)\;\forall n\in\Z_{\geq1}\;\exists u\in W(F):\quad u^n=h.
\]
\textbf{Unresolved qualification.} The supplied wording does not require the subgroup to be nontrivial. The trivial group satisfies the definition of radicability, so the intended substantive question requires clarification about whether \(H\ne\{1\}\), respectively \(W(F)\ne\{1\}\), is required. These restrictions have not been silently added. The source also specifies no variety, rank, finiteness condition, or properness requirement; the unrestricted quantifiers above record its literal scope, subject to this unresolved qualification.
\subsection{Short English statement}
Brumberg asks whether a relatively free group can contain a radicable subgroup, and more specifically a radicable verbal subgroup. Radicability requires every positive integer root to exist within the subgroup. The supplied statement leaves nontriviality unstated, so its intended substantive formulation needs human review.
\subsection{Sources}
\begin{itemize}
\item Local UnsolvedMath registry, permanent ID \(700159\), source problem code \texttt{catalogue entry}, title \emph{Kourovka Notebook Problem 3.5}; supplied upstream revision \texttt{d371ee917ac9647fbb06cd6fc529deec6f3dda82}.
\item The supplied record attributes Problem 3.5 to N. R. Brumberg, proposed in 1969, and identifies its source as \emph{The Kourovka Notebook (Unsolved Problems in Group Theory)}, 3rd issue (1969), as included in the 21st issue edition (2026), \texttt{arXiv:1401.0300v48}.
\item The supplied collection description names E. I. Khukhro and V. D. Mazurov as editors of the Kourovka Notebook collection covering Issues 1--21 (1965--2026).
\item The local record labels the problem \texttt{open} and records an update timestamp of \texttt{2026-10-16T14:00:00Z}. These are supplied registry metadata, not an independently verified assessment of the literature. The cited edition has not been independently consulted for this draft.
\end{itemize}
\end{document}
